\clearpage\label{answers}
{\sans\bfseries\LARGE\color{ink} Answers}\par
\footnotesize
Try every puzzle first. Where an answer shows a calculation, check the steps, not just the result.

\subhead{Chapter 1}
\textbf{Multix on stillness.} No. A book on a table has two arrows, gravity down and the table's push up, and they pop to nothing. \textbf{Keep box.} Kept: the direction and the size of each push. Set aside: the exact shape of the thing.

\textbf{1.1} Gravity pulls the bag down. The string pulls the bag up.\\
\textbf{1.2} (a) $6-2=4$ right. (b) 0. (c) $5-3=2$ left.\\
\textbf{1.3} Answers depend on the diagram; check that the leftover points the way the thing moves.\\
\textbf{1.4} Gravity down, the floor's push up, same size.\\
\textbf{1.5} Push or pull, by your own check: a door handle is a pull to open it, a hand on a wall is a push.\\
\textbf{1.6} Another arrow of size 2 to the right. \textbf{1.7} 1 to the right. \textbf{1.8} $7-4=3$ to the right, so the boat goes right. \textbf{1.9} Answers vary.

\subhead{Chapter 2}
\textbf{Multix on friction.} No. The stretch of the rubber band is a stand-in for the pull, and 3~cm is a measure of the pull, not a force in other units. \textbf{Keep box.} Kept: friction as one arrow against sliding. Set aside: every bump.

\textbf{2.1} Rug: $22+25+23=70$, average $70\div3\approx23$~cm. Floor: $62+58+60=180$, average 60~cm.\\
\textbf{2.2} Carpet stops it soonest. Its slide distance is the smallest.\\
\textbf{2.3} Look for two changes at once, such as a different push and a different surface.\\
\textbf{2.4} Friction helps shoes grip, brakes stop wheels, and knots hold.\\
\textbf{2.5} A steady pull means the pull and friction are equal.\\
\textbf{2.6} Card has more friction. $3.1-1.9=1.2$~cm. \textbf{2.7} About 12~cm (11 to 13 is fine). It is a guess because we did not test four blocks. \textbf{2.8} Left, then right: friction points against the sliding. \textbf{2.9} Answers vary.

\subhead{Chapter 3}
\textbf{Multix on falling.} The two do not fall together because the air pushes up more on the flat paper. \textbf{Keep box.} Kept: gravity down and an air push up. Set aside: the swirls in the air.

\textbf{3.1} Crumpled: $0.6+0.7+0.6+0.8+0.6=3.3$, average $3.3\div5=0.66\approx0.7$~s. Flat: $1.9+2.2+2.0+2.1+1.8=10.0$, average 2.0~s.\\
\textbf{3.2} Gravity down, air push up, nearly equal.\\
\textbf{3.3} Gravity down, air push up, same size in the middle of the fall. A heavy flat boat can float because a wide boat gets pushed up by more water.\\
\textbf{3.4} Marco changed the paper and the drop height. He should change only one.\\
\textbf{3.5} $2.0-0.7=1.3$~s. About 3 times as long.\\
\textbf{3.6} The fall time gets longer, because the air's push up is bigger and the leftover is smaller. \textbf{3.7} 2 times. A wider boat pushes aside more water, so the water pushes up more. \textbf{3.8} $5-1=4$ down. Later the air's push is 5.

\subhead{Chapter 4}
\textbf{Multix on the seesaw.} No. $3\times4=12$ and $1\times4=4$. They are not equal, so the left side goes down. \textbf{Keep box.} Kept: coins times distance. Set aside: the weight of the ruler.

\textbf{4.1} $2\times9=18$, and $3\times6=18$, so 6~cm.\\
\textbf{4.2} Left: $3\times5=15$. Right: $2\times8=16$. The right side goes down.\\
\textbf{4.3} $5\times n=2\times10=20$, so $n=4$ coins.\\
\textbf{4.4} $30\times2=60$ and $40\times1.5=60$, so 1.5~m.\\
\textbf{4.6} $3\times10=30$ and $30\div2=15$~cm. \textbf{4.7} $5\times6=30$, so 30~cm, which is past the end of a 30~cm ruler. Use two coins at 15~cm. \textbf{4.8} $4\times3=12$ and $3\times4=12$. Yes. \textbf{4.9} $3\times8=24$ and $24\div6=4$ coins.\\
\textbf{4.5} $2+6=8$ and $1+12=13$, so the rule ``plus'' does not work. Coins times distance does.

\subhead{Chapter 5}
\textbf{Multix on the ramp.} No. Gravity is the same. A steeper ramp lets more of the pull act along the slope. \textbf{Keep box.} Kept: the directions of the forces. Set aside: the shape of the board.

\textbf{5.1} $55+57+59=171$, average $171\div3=57$~cm.\\
\textbf{5.2} Any guess above 132 is fine, for example 160. It is only a guess, because we did not test four books.\\
\textbf{5.3} The 4~m ramp needs a smaller push, but you push it twice as far.\\
\textbf{5.4} Two things changed, the ramp and the floor.\\
\textbf{5.5} Gravity down, the slope's push out from the slope, and friction up the slope. Nothing is missing.\\
\textbf{5.6} Friction.\\
\textbf{5.7} $100-62=38$ and $132-100=32$. \textbf{5.8} $6.0-3.8=2.2$ and $6.0-2.5=3.5$. \textbf{5.9} The steep one; the arrow along the slope is bigger. \textbf{5.10} Answers vary.

\subhead{Chapter 6}
\textbf{Multix on the sled.} No. The steady sled has no leftover. Friction is small, so the arrows pop to nearly nothing. \textbf{Keep box.} Kept: direction and size. Set aside: which arrow is which.

\textbf{6.1} Friction, pointing left, size 4.\\
\textbf{6.2} The table's push, up, size 3.\\
\textbf{6.3} (a) $7-3-2=2$ right. (b) 0. (c) $9-4=5$ left.\\
\textbf{6.4} Left: $3\times4=12$. Right: $2\times5=10$. Left wins by 2.\\
\textbf{6.5} Water pushes up with size 10.\\
\textbf{6.6} Carpet. It has the most friction.\\
\textbf{6.7} $4\times5=20$ and $2\times10=20$, so 2 coins. Check: if the bag is at 10~cm and the coins at 5~cm, the bag is 2 coins.\\
\textbf{6.8} None. The ice pushes up and gravity pulls down.\\
\textbf{6.9} Air pushes up with size 2.\\
\textbf{6.10} 0. \textbf{6.11} It speeds up. \textbf{6.12} $2\times12=24$ and $4\times6=24$. Yes. \textbf{6.13} Check that the arrows explain what the thing did. \textbf{6.14} Gravity down 2, air up 2. \textbf{6.15} $7-3=4$ right. \textbf{6.16} Gravity down and water up, equal.\\
\textbf{Slop check.} (a) You cannot tell: no title, no labels. (b) Many colours, stars, joke text, tilted filler. (c) A title, the surfaces and the distances with units. (d) The stars and jokes.

\clearpage\label{notes}
{\sans\bfseries\LARGE\color{ink} Words}\par
\small
\begin{multicols}{2}
\word{balance}{when the turning effects on both sides are equal.}
\word{force}{a push or a pull.}
\word{friction}{a force that points against sliding.}
\word{gravity}{the pull that draws things toward the Earth.}
\word{leftover}{what is left when the arrows are popped.}
\word{model}{a simple stand-in that keeps some things and drops others.}
\word{pivot}{the point a seesaw balances on.}
\word{ramp}{a slope used to lift things gently.}
\word{spread}{how far apart repeated readings are.}
\end{multicols}

\vspace{0.4em}
{\sans\bfseries\large\color{sky!70!black} Notes for grown-ups}\par
This book teaches that a force diagram is a model: it reduces a push or pull to an arrow with a direction and a size. Each chapter has one activity, one worked example with made-up practice numbers, one trust box, one ``what did we keep?'' prompt and a page of puzzles.

\textbf{Materials and safety.} A ruler, a pencil, ten coins, rubber bands, a board, books, a toy car and string. Keep ramps low and away from stairs. Do not drop things from windows or stairs. Rubber bands can snap: do not stretch them toward faces.

\textbf{The practice numbers} were made up so the arithmetic is easy to check. Real results depend on the surfaces and the objects, so real results will differ.

\subhead{Sources}
\begin{itemize}[leftmargin=1.3em,itemsep=0.2em]
\item J. R. Taylor. \emph{An Introduction to Error Analysis}, 2nd edition. University Science Books, 1997. For repeats and spread.
\item D. Halliday, R. Resnick and J. Walker. \emph{Fundamentals of Physics}. Wiley. For forces, friction and balance.
\end{itemize}

\clearpage
\thispagestyle{blank}
\begin{tikzpicture}[remember picture,overlay]
  \fill[sun!35] ($(current page.south west)+(0.4in,0.4in)$) rectangle ($(current page.north east)-(0.4in,0.4in)$);
  \node[anchor=north,align=center,font=\sffamily\bfseries,text=ink,text width=5.4in] at ($(current page.north)+(0,-1.4in)$) {\Large Next in Science\\[6pt]\fontsize{28}{30}\selectfont Where the Water Goes};
  \node[align=center,font=\sffamily,text width=4.6in,text=ink] at (current page.center) {A cup of water sits on a shelf, and a day later there is less. Book 5 follows water as it changes and moves, and keeps track of how much there was.};
  \node at ($(current page.south)+(0,2.6in)$) {\scalebox{3}{\multixbot[ok]}};
  \node[font=\sffamily\bfseries,text=ink] at ($(current page.south)+(0,0.9in)$) {Ages 9 to 11 \quad$\star$\quad Science Book 4};
\end{tikzpicture}
